En una actividad de topografía escolar, un estudiante ubica tres puntos en el plano que representan la ubicación de tres señales en el terreno. Al intentar alinearlos con una regla, observa que no están sobre una misma línea recta.

Sin embargo, el estudiante afirma que es posible trazar una única recta que pase por los tres puntos si se ajusta correctamente la inclinación.

¿La afirmación es verdadera o falsa?

Justifique su respuesta utilizando argumentos geométricos claros y explicando la condición necesaria para que varios puntos sean colineales. No se aceptan respuestas sin justificación.